\subsection{2020 Census: Webster vs.\ Huntington-Hill}

We compute the Webster apportionment for 2020 Census data using
\textsc{bisect-apportion}'s Webster implementation.
The Webster divisor is found by binary search starting from the natural
divisor $d_0 = 761{,}169$ and adjusting until
$\sum_i \max(1, \lfloor p_i/d + 0.5 \rfloor) = 435$.

\textbf{Result}: For the 2020 Census, Webster produces the same seat
allocation as Huntington-Hill for all 50 states.

This agreement is not coincidental: it reflects the fact that for the
2020 Census population distribution, no state's exact quota (using the
Webster divisor) falls in the interval $(\sqrt{s(s+1)}, s+0.5)$ for
any $s$.
In other words, all states are either clearly above or clearly below the
rounding threshold under both methods.

\subsection{States Near the Boundary}

The states closest to a rounding threshold under Webster in 2020 are
those whose exact quotas $q_i = p_i \times 435 / 331{,}108{,}434$
have fractional parts nearest to 0.5.

\begin{table}[h]
\centering
\caption{States with exact quotas nearest to 0.5 rounding thresholds,
  2020 Census (Webster method).}
\label{tab:near-boundary}
\begin{tabular}{lrrrr}
\toprule
State & Population & Exact quota & Fractional part & Webster seats \\
\midrule
Montana   & 1,084,225  & 1.425 & 0.425 & 2 \\
Rhode Island & 1,097,379 & 1.442 & 0.442 & 2 \\
New Hampshire & 1,377,529 & 1.811 & 0.811 & 2 \\
Hawaii    & 1,455,271  & 1.912 & 0.912 & 2 \\
\bottomrule
\end{tabular}
\begin{tablenotes}
\small
\item States with fractional parts far from 0.5 are not affected by the
  HH vs.\ Webster choice. Montana and Rhode Island, with fractional parts
  near 0.42--0.44, are the states closest to a different outcome under
  a small perturbation of the divisor.
\end{tablenotes}
\end{table}

Montana is the most historically significant near-threshold state:
its 2020 population of 1{,}084{,}225 gives it a quota of $\approx 1.425$
under both HH and Webster.
Under both methods, this rounds to 2 seats.
Under the 1990 Census, Montana's lower population placed it below the
HH threshold for a second seat (receiving 1), triggering the
\textit{Montana} litigation.

\subsection{Historical Census Comparisons}

Webster and Huntington-Hill do not always agree.
For the 1940 Census (the last apportionment under Webster before
the 1941 Act adopted HH), the methods differed for Michigan:
Webster would give Michigan 17 seats, HH gives 17 seats for that
census --- the methods happened to agree.
More interesting is the comparison on hypothetical House sizes:
at $H = 400$, Webster and HH differ by 1 seat for at least one state
in the 2020 Census distribution; at $H = 435$, they agree.
This size-dependence illustrates that the methods' agreement or
disagreement is a function of the specific population distribution
and House size, not a general rule.
